\clearpage
% Entry point for the new two-step bridge. Existing C053 diagrams are untouched.
\begin{proofbox}
\textbf{\textcolor{CProof}{方法一的思路：第一步，先只看初中几何}}

在直角梯形 $MABN$ 中，$AM\parallel BN$。$P$ 是斜腰 $AB$ 上的普通点，
过 $P$ 作平行于底边的 $PQ$，并作垂直于底边的 $HK$。
记 $AP=r$、$PB=s$、$PQ=h$。
\begin{center}% Pure Euclidean geometry only: no conic, focus, directrix or coordinates.
% Same actual points and scale as C053_trapezoid_parabola_mapping.tex.
\begin{tikzpicture}[font=\small,>=stealth]
\draw[rounded corners=3pt,gray!45] (0,0) rectangle (14.2,11.5);
\node at (7.1,10.95) {先看一个直角梯形：上底怎样加，下底怎样减？};
\begin{scope}[shift={(1.6,5)},scale=1.15]
\fill[blue!7] (2,0)--(2,3.464101615)--(4,3.464101615)--cycle;
\fill[orange!10] (2,0)--(2,-1.154700538)--(1.333333333,-1.154700538)--cycle;
\draw[gray!80,thick] (0,3.464101615)--(0,-1.154700538);
\draw[blue!65!black,thick] (0,3.464101615)--(4,3.464101615);
\draw[orange!85!black,thick] (0,-1.154700538)--(1.333333333,-1.154700538);
\draw[blue!75!black,very thick] (2,0)--(4,3.464101615)
 node[midway,below right=5pt] {$r$};
\draw[orange!85!black,very thick] (2,0)--(1.333333333,-1.154700538)
 node[midway,above left=5pt] {$s$};
\draw[gray,dashed] (2,-1.154700538)--(2,3.464101615);
\draw[gray,dashed] (0,0)--(3.15,0);
\node[above=5pt] at (.65,0) {$h$};
\draw[blue!65!black,densely dashed] (1.333333333,-1.154700538)--(2,-1.154700538);
\draw[gray!80] (0,3.26)--(.2,3.26)--(.2,3.464101615);
\draw[gray!80] (0,-.95)--(.2,-.95)--(.2,-1.154700538);
\draw[blue!65!black] (2,3.26)--(2.2,3.26)--(2.2,3.464101615);
\draw[blue!65!black] (1.8,-1.154700538)--(1.8,-.954700538)--(2,-.954700538);
\draw[violet,->] (2.55,0) arc[start angle=0,end angle=60,radius=.55];
\node[violet] at (2.83,.34) {$\theta$};
\fill (0,3.464101615) circle(2pt) node[left=5pt] {$M$};
\fill (4,3.464101615) circle(2pt) node[above=8pt] {$A$};
\fill (1.333333333,-1.154700538) circle(2pt) node[below left=6pt] {$B$};
\fill (0,-1.154700538) circle(2pt) node[left=5pt] {$N$};
\fill (0,0) circle(1.5pt) node[left=5pt] {$Q$};
\fill (2,0) circle(2pt) node[above left=6pt] {$P$};
\fill (2,3.464101615) circle(1.5pt) node[above left=5pt] {$H$};
\fill (2,-1.154700538) circle(1.5pt) node[right=5pt] {$K$};
\node[above=6pt,blue!65!black] at (3,3.464101615) {$r\cos\theta$};
\draw[gray,dotted] (1.333333333,-1.55)--(1.333333333,-2.05);
\draw[gray,dotted] (2,-1.55)--(2,-2.05);
\draw[<->,orange!85!black] (1.333333333,-2)--(2,-2)
 node[midway,below] {$s\cos\theta$};
\end{scope}
\node[anchor=west,align=left] at (8.6,7)
 {$AM\parallel QP\parallel BN$\\[10pt]
 $AP=r,\quad PB=s$\\[9pt]$PQ=h$\\[13pt]
 $MH=NK=PQ=h$\\[10pt]
 $HA=r\cos\theta$\\[9pt]$BK=s\cos\theta$};
\node at (7.1,1.5) {看平行线和直角三角形：$AM=MH+HA$，$BN=NK-BK$。};
\node at (7.1,.6) {$\displaystyle AM=h+r\cos\theta,\qquad BN=h-s\cos\theta$};
\end{tikzpicture}
\end{center}
图示 $0<\theta<90^\circ$。在直角三角形 $PHA$、$PKB$ 中，水平直角边分别是
$r\cos\theta$、$s\cos\theta$；又 $MH=NK=PQ=h$，于是上底相加、下底相减。
\textbf{先记住这个梯形的长度分解，再给图中的点和线赋予新的含义。}
\end{proofbox}
\clearpage
\begin{proofbox}
\textbf{\textcolor{CProof}{方法一的思路：第二步，把同一个梯形放回抛物线}}

梯形的大小、形状和辅助线都与上一图相同。
现在令 $P$ 对应焦点 $F$，$MN$ 所在直线对应准线，故 $h=FQ=p$。
\begin{center}% Same trapezoid and coordinates as the pure geometry figure, translated
% from y^2=4x by x -> x+1. Here F=(2,0), directrix x=0 in drawing coordinates.
\begin{tikzpicture}[font=\small,>=stealth]
\draw[rounded corners=3pt,gray!45] (0,0) rectangle (14.2,11.5);
\node at (7.1,10.95) {再放回抛物线：只补上 $AM=AF$、$BN=BF$ 两个等量};
\begin{scope}[shift={(1.6,5)},scale=1.15]
\fill[blue!7] (2,0)--(2,3.464101615)--(4,3.464101615)--cycle;
\fill[orange!10] (2,0)--(2,-1.154700538)--(1.333333333,-1.154700538)--cycle;
\draw[teal!75!black,thick,samples=120,domain=-2.3:4.15,variable=\t]
 plot ({1+\t*\t/4},{\t});
\draw[violet,densely dashed] (0,-2.35)--(0,4.2);
\node[violet,rotate=90] at (-.65,1.5) {准线};
\draw[gray!80,thick] (0,3.464101615)--(0,-1.154700538);
\draw[blue!65!black,thick] (0,3.464101615)--(4,3.464101615);
\draw[orange!85!black,thick] (0,-1.154700538)--(1.333333333,-1.154700538);
\draw[blue!75!black,very thick] (2,0)--(4,3.464101615)
 node[midway,below right=5pt] {$r$};
\draw[orange!85!black,very thick] (2,0)--(1.333333333,-1.154700538)
 node[midway,above left=5pt] {$s$};
\draw[gray,dashed] (2,-1.154700538)--(2,3.464101615);
\draw[gray,dashed] (0,0)--(3.15,0);
\node[above=5pt] at (.65,0) {$p$};
\draw[blue!65!black,densely dashed] (1.333333333,-1.154700538)--(2,-1.154700538);
\draw[gray!80] (0,3.26)--(.2,3.26)--(.2,3.464101615);
\draw[gray!80] (0,-.95)--(.2,-.95)--(.2,-1.154700538);
\draw[blue!65!black] (2,3.26)--(2.2,3.26)--(2.2,3.464101615);
\draw[blue!65!black] (1.8,-1.154700538)--(1.8,-.954700538)--(2,-.954700538);
\draw[violet,->] (2.55,0) arc[start angle=0,end angle=60,radius=.55];
\node[violet] at (2.83,.34) {$\theta$};
\fill (0,3.464101615) circle(2pt) node[left=5pt] {$M$};
\fill (4,3.464101615) circle(2pt) node[above=8pt] {$A$};
\fill (1.333333333,-1.154700538) circle(2pt) node[below left=6pt] {$B$};
\fill (0,-1.154700538) circle(2pt) node[left=5pt] {$N$};
\fill (0,0) circle(1.5pt) node[left=5pt] {$Q$};
\fill (2,0) circle(2pt) node[above left=6pt] {$F$};
\fill (2,3.464101615) circle(1.5pt) node[above left=5pt] {$H$};
\fill (2,-1.154700538) circle(1.5pt) node[right=5pt] {$K$};
\end{scope}
\node[anchor=west,align=left] at (8.6,7)
 {上一图：$P\longrightarrow F$\\[9pt]$h\longrightarrow p$\\[14pt]
 抛物线的等距定义：\\[10pt]
 $\textcolor{blue!75!black}{AM=AF=r}$\\[10pt]
 $\textcolor{orange!85!black}{BN=BF=s}$\\[14pt]
 梯形的分解仍然成立。};
\node at (7.1,1.5) {几何关系给“加减”，抛物线定义给“相等”。};
\node at (7.1,.6) {$\displaystyle r=p+r\cos\theta,\qquad s=p-s\cos\theta$};
\end{tikzpicture}
\end{center}
抛物线定义只补上两件事：$AM=AF=r$、$BN=BF=s$。
把它们代入上一图的两条几何关系，就得到图下的两个方程。因此
\[
 r=\frac p{1-\cos\theta},\quad s=\frac p{1+\cos\theta},\qquad
 L=r+s=\frac{2p}{\sin^2\theta}.
\]
两图取 $p=2$、$\theta=60^\circ$；一般方向按有向投影写方程。
\textbf{下面保留原来的两张逐端投影图，分别对照 $AF$ 和 $BF$。}
\end{proofbox}

\clearpage
\begin{proofbox}
\textbf{\textcolor{CProof}{方法一：焦点准线定义＋两端投影（先求正向一段）}}

取右开口 $y^2=2px$，$F=(p/2,0)$，准线 $x=-p/2$。
沿 $(\cos\theta,\sin\theta)$ 到 $A$，反向到 $B$，令 $r=FA$、$s=FB$。
\begin{center}\begin{tikzpicture}[font=\small,>=stealth]
\draw[rounded corners=3pt,gray!45] (0,0) rectangle (14.2,12);
\node at (7.1,11.45) {先看正方向：到准线的距离 $AM=p+r\cos\theta$};
\begin{scope}[shift={(3.8,5.6)},scale=1.1]
\draw[->,gray] (-1.5,0)--(4.5,0) node[above] {$x$};

\draw[violet,densely dashed] (-1,-2.3)--(-1,4.2);
\node[rotate=90,violet] at (-1.4,1.5) {准线};
\fill[blue!7] (1,0)--(1,3.464101615)--(3,3.464101615)--cycle;
\draw[teal!75!black,thick,samples=120,domain=-2.3:4.15,variable=\t] plot ({\t*\t/4},{\t});
\draw[gray,dashed] (1,0)--(1,3.464101615);
\draw[orange!85!black,thick,dashed] (-1,3.464101615)--(1,3.464101615);
\draw[blue!75!black,thick,dashed] (1,3.464101615)--(3,3.464101615)
 node[midway,above=5pt] {$r\cos\theta$};
\draw[blue!75!black,very thick] (1,0)--(3,3.464101615) node[midway,below right=5pt] {$r$};
\fill (3,3.464101615) circle(2pt) node[above=8pt] {$A$};
\fill (-1,3.464101615) circle(1.5pt) node[left=4pt] {$M$};
\fill (1,3.464101615) circle(1.5pt) node[below left=5pt] {$H$};
\draw[gray,dotted] (1,3.8)--(1,4.2);
\draw[<->,orange!85!black] (-1,4.15)--(1,4.15) node[midway,above] {$MH=p$};
\draw[blue!75!black] (1,3.25)--(1.21,3.25)--(1.21,3.464101615);
\draw[violet,->] (1.65,0) arc[start angle=0,end angle=60,radius=.65];
\node[violet] at (1.94,.42) {$\theta$};
\fill (1,0) circle(2pt) node[below right=5pt] {$F$};

\end{scope}
\node[anchor=west,align=left] at (9.3,7)
 {$\textcolor{blue!75!black}{FA=r}$\\[10pt]$AM=FA=r$\\[12pt]
 $MH=p$\\[9pt]$HA=r\cos\theta$\\[12pt]$AM=MH+HA$};
\node at (7.1,1.6) {同一模型：$y^2=4x$，$\theta=60^\circ$；蓝色三角形给水平投影。};
\node at (7.1,.65) {$\displaystyle r=p+r\cos\theta\quad\Longrightarrow\quad r=\frac{p}{1-\cos\theta}$};
\end{tikzpicture}
\end{center}
\textbf{先看 $AF$：}图中 $H$ 在 $F$ 的正上方，$MH=p$，$HA=r\cos\theta$。
由抛物线定义 $FA=AM$，一般地按有向投影写成
\[
 r=p+r\cos\theta\quad\Longrightarrow\quad
 \boxed{r=\frac{p}{1-\cos\theta}}.
\]
两张投影图取 $p=2,\theta=60^\circ$ 便于观察；公式用有向分量，对任意 $0<\theta<\pi$ 均成立。
\end{proofbox}
\clearpage
\begin{proofbox}
\textbf{\textcolor{CProof}{方法一（续）：补齐反向一段，再把两段相加}}
\begin{center}\begin{tikzpicture}[font=\small,>=stealth]
\draw[rounded corners=3pt,gray!45] (0,0) rectangle (14.2,12);
\node at (7.1,11.45) {再看反方向：从 $p$ 中减去水平投影，$BN=p-s\cos\theta$};
\begin{scope}[shift={(3.8,5.6)},scale=1.1]
\draw[->,gray] (-1.5,0)--(4.5,0) node[above] {$x$};

\draw[violet,densely dashed] (-1,-3.2)--(-1,4.15);
\node[rotate=90,violet] at (-1.4,1.5) {准线};
\fill[orange!10] (1,0)--(1,-1.154700538)--(.333333333,-1.154700538)--cycle;
\draw[teal!75!black,thick,samples=120,domain=-1.65:4.15,variable=\t] plot ({\t*\t/4},{\t});
\draw[gray,dashed] (1,0)--(1,-1.154700538);
\draw[orange!85!black,thick,dashed] (-1,-1.154700538)--(.333333333,-1.154700538);
\draw[blue!65!black,dashed] (.333333333,-1.154700538)--(1,-1.154700538);
\draw[orange!85!black,very thick] (1,0)--(.333333333,-1.154700538);
\fill (.333333333,-1.154700538) circle(2pt) node[below left=5pt] {$B$};
\fill (-1,-1.154700538) circle(1.5pt) node[left=4pt] {$N$};
\fill (1,-1.154700538) circle(1.5pt) node[right=5pt] {$K$};
\draw[blue!65!black] (.82,-1.154700538)--(.82,-.974700538)--(1,-.974700538);
\draw[gray,dotted] (.333333333,-1.6)--(.333333333,-2.1);
\draw[gray,dotted] (1,-1.55)--(1,-3.05);
\draw[gray,dotted] (-1,-1.5)--(-1,-3.05);
\draw[<->,blue!65!black] (.333333333,-2.05)--(1,-2.05) node[midway,below] {$BK=s\cos\theta$};
\draw[<->,violet] (-1,-3)--(1,-3) node[midway,below] {$NK=p$};
\fill (1,0) circle(2pt) node[below right=5pt] {$F$};

\end{scope}
\node[anchor=west,align=left] at (9.3,7)
 {$\textcolor{orange!85!black}{FB=s}$\\[10pt]$BN=FB=s$\\[12pt]
 $NK=p$\\[9pt]$BK=s\cos\theta$\\[12pt]$BN=NK-BK$};
\node at (7.1,1.25) {与上一图对照：同一抛物线、同一方向、同一尺度。};
\node at (7.1,.5) {$\displaystyle s=p-s\cos\theta\quad\Longrightarrow\quad s=\frac{p}{1+\cos\theta}$};
\end{tikzpicture}
\end{center}
\textbf{再看 $BF$：}图中 $K$ 在 $F$ 的正下方，$NK=p$，$BK=s\cos\theta$。
到准线的距离是 $BN=NK-BK$，由 $FB=BN$ 得
\[
 s=p-s\cos\theta\quad\Longrightarrow\quad
 \boxed{s=\frac{p}{1+\cos\theta}}.
\]
两分母在 $0<\theta<\pi$ 内均为正；焦点处于两个交点之间，故
\[
 L=r+s=\frac p{1-\cos\theta}+\frac p{1+\cos\theta}
 =\frac{2p}{1-\cos^2\theta}=\boxed{\frac{2p}{\sin^2\theta}}.
\]
\textbf{记忆：}正向减余弦，反向加余弦；两段相加，平方差给出正弦平方。
\end{proofbox}
\clearpage
\begin{proofbox}
\textbf{\textcolor{CProof}{方法二（门槛最低，推荐不熟悉三角投影时）：只求一元方程的“根之和”}}\par
\textbf{优势：一元二次方程，不解根、不求判别式、不算坐标距离。}
仍以 $y^2=2px$ 为例：
\begin{enumerate}[leftmargin=*,itemsep=.35em]
\item \textbf{先把弦长改为横坐标的和。} 抛物线定义给出
$|AF|=x_A+\frac p2$、$|BF|=x_B+\frac p2$，焦点在两交点之间，故
\[\boxed{L=x_A+x_B+p}.\]
\item \textbf{只设一个未知量 $y$。} 设过焦点的直线为
\[x=\frac p2+y\cot\theta.\]
这是把 $\Delta x/\Delta y=\cot\theta$ 写成方程；即使竖直 $\theta=\pi/2$ 也成立。
\item \textbf{代入抛物线，只用韦达定理“根的和”。}
\[
y^2=2p\left(\frac p2+y\cot\theta\right)
\iff y^2-2p\cot\theta\,y-p^2=0.
\]
其两根是 $y_A,y_B$，因此
\[y_A+y_B=2p\cot\theta.\]
又 $x_A+x_B=p+(y_A+y_B)\cot\theta=p+2p\cot^2\theta$。
\item \textbf{最后一加。}
\[
L=x_A+x_B+p=2p(1+\cot^2\theta)
=\boxed{\frac{2p}{\sin^2\theta}}.
\]
\end{enumerate}
\textbf{一行复盘：焦半径求和 $\to$ 只求 $y_A+y_B$ $\to$ $1+\cot^2=\csc^2$。}
\end{proofbox}
\clearpage
\begin{proofbox}
\textbf{\textcolor{CProof}{方法三（常规标准证明）：单位方向参数方程 + 韦达定理}}\par
设 $F=(p/2,0)$，焦点弦所在直线上点的参数方程为
\[
x=\frac p2+t\cos\theta,\qquad y=t\sin\theta.
\]
方向向量 $(\cos\theta,\sin\theta)$ 长度为 $1$，故参数 $t$ 是\textbf{有向距离}，
两个交点对应 $t_1>0>t_2$，与前图一致：$t_1=r$、$t_2=-s$，故 $L=t_1-t_2=r+s$。
代入 $y^2=2px$：
\[
t^2\sin^2\theta-2pt\cos\theta-p^2=0,
\]
于是
\[
t_1+t_2=\frac{2p\cos\theta}{\sin^2\theta},\qquad
t_1t_2=-\frac{p^2}{\sin^2\theta}.
\]
用恒等式 $(t_1-t_2)^2=(t_1+t_2)^2-4t_1t_2$：
\[
\begin{aligned}
L^2&=\frac{4p^2\cos^2\theta}{\sin^4\theta}
+\frac{4p^2}{\sin^2\theta}\\
&=\frac{4p^2(\cos^2\theta+\sin^2\theta)}{\sin^4\theta}
=\frac{4p^2}{\sin^4\theta}.
\end{aligned}
\]
因 $L>0$，故 $\boxed{L=2p/\sin^2\theta}$。
\end{proofbox}
\clearpage
\begin{proofbox}
\textbf{\textcolor{CProof}{方法四（便于笔算的斜率法，供比较）：$k$ 法}}\par
若直线有非零有限斜率 $k$，则
$y=k(x-\frac p2)$，即 $x=\frac p2+\frac yk$。
代入抛物线得到 $y^2-\frac{2p}{k}y-p^2=0$，
由韦达定理 $y_A+y_B=\frac{2p}k$，所以
\[
L=x_A+x_B+p
=2p+\frac{y_A+y_B}k
=\boxed{\frac{2p(1+k^2)}{k^2}}.
\]
\textbf{适用提醒：}$k=0$ 时不构成两交点的焦点弦；竖直直线斜率不存在，改用 $L=2p$。
\medskip
\begin{center}\begin{tikzpicture}[font=\small,>=stealth]
\draw[rounded corners=3pt,gray!45] (0,0) rectangle (14.2,10.7);
\node at (7.1,10.149999999999999) {左右开口互为镜像：焦点换边，弦长保持不变};
\node at (3.6,9.3) {右开口 $y^2=4x$};
\begin{scope}[shift={(2.4,5.2)},scale=.85]
\draw[->,gray] (-1.1,0)--(4.5,0);
\draw[teal!75!black,thick,samples=120,domain=-2.2:4,variable=\t] plot ({1*\t*\t/4},{\t});
\draw[blue!75!black,very thick] (1,0)--(3,3.464101615);
\draw[orange!85!black,very thick] (1,0)--(0.3333333333333333,-1.154700538);
\fill (1,0) circle(2pt) node[below right=4pt] {$F$};
\fill (3,3.464101615) circle(2pt) node[above=8pt] {$A$};
\fill (0.3333333333333333,-1.154700538) circle(2pt) node[below left=6pt] {$B$};
\end{scope}
\node at (3.6,2.3) {$\theta=60^\circ,\quad L=16/3$};
\node at (10.6,9.3) {左开口 $y^2=-4x$};
\begin{scope}[shift={(11.8,5.2)},scale=.85]
\draw[->,gray] (-4.5,0)--(1.1,0);
\draw[teal!75!black,thick,samples=120,domain=-2.2:4,variable=\t] plot ({-1*\t*\t/4},{\t});
\draw[blue!75!black,very thick] (-1,0)--(-3,3.464101615);
\draw[orange!85!black,very thick] (-1,0)--(-0.3333333333333333,-1.154700538);
\fill (-1,0) circle(2pt) node[below left=4pt] {$F$};
\fill (-3,3.464101615) circle(2pt) node[above=8pt] {$A$};
\fill (-0.3333333333333333,-1.154700538) circle(2pt) node[below right=6pt] {$B$};
\end{scope}
\node at (10.6,2.3) {$\theta=120^\circ,\quad L=16/3$};

\node at (7.1,1.2) {$\sin^2(180^\circ-\theta)=\sin^2\theta$，镜像前后长度相同。};
\node at (7.1,.5) {两边都取 $p=2>0$；开口向左，不代表长度带负号。};
\end{tikzpicture}
\end{center}
左右开口的图形关于 $y$ 轴镜像；长度不变，倾斜角变为 $\pi-\theta$，正弦平方不变。
因此左右开口均有 $L=2p/\sin^2\theta$，$p$ 都取正数。
\par\medskip
\textbf{\textcolor{CProof}{通径最短的证明}}\par
$0<\sin^2\theta\le1$，故 $L\ge2p$；
等号 $\sin^2\theta=1\iff \theta=\frac\pi2$，此时弦与抛物线对称轴垂直。
\end{proofbox}
